\section*{Required author factorization lemmas, Sections 5.3.2--5.3.3}
\paragraph{Lemma 5.23 ([5],Lemmas6.14/7.14).} Suppose $(R_l,R_m,R_r)\to(R'_l,R^*_m,R'_r)$. In Definition5.21's notation: (1) $|S'_l|+|S'_r|\ge|A|+|B|+1$. (2) Put $s=(|A|+|B|)/2$, let $p'$ be the maximum number of vertex-disjoint paths from $S'_l$ to $S'_r$ in $R^*_m$, and $p$ from $A$ to $B$ in $R_m$. Then
\[2(s'-s)+(p-p')+|I(R^*_m)|\le z(R_l,R_m,R_r).\]
\begin{proof}(1) There must be an unexpected pairwise intersection. In any of the three cases breaking (5.26), there is $v\in Z$ in $V(R_l)-A$ or $V(R_r)-B$. WLOG the first. Then $S'_l\ne A$ since by left-generatedness $v$ is reachable from $I$ without passing $A$. Similarly, if $|S'_l|=|A|$ it is $A$, since $A$ is the unique min-separator separating $(I,A)$, impossible. Thus $S'_l>A$. (2) The author refers to the proof in the cited original paper.\end{proof}
\footnote{In this setting $R_m$ is a ribbon without isolated vertices. The credited prior lemma is a standard prerequisite, not a separately selected contribution.}
The separating factorization cancels a term in $[L(DQ_0D)L^\top]_{\mathrm{non\text{-}can}}$ by a canonical term using $R'_l$, $R'_r$ and the largest ribbon of $R^*_m$ in a new middle matrix $DQ_1D$. Other triples can have the same separating factorization, and entries of $L$ are sums of different left ribbons. Simultaneous cancellation requires the following refined solvability condition, also needed in Lemma6.20.
\paragraph{Proposition 5.24 (Solvability condition, cf.[5],Claim6.12).} Fix $(I,J,S'_l,S'_r)$ and generalized ribbon $R^*_m$ on $(S'_l,S'_r)$. Let $(R'_l,R'_r)$ be inner-canonical left/right ribbons on $(I,S'_l)$, $(S'_r,J)$. Let $(R''_l,R''_r)$ be another such pair with $e(R'_l)=e(R''_l)$, $e(R'_r)=e(R''_r)$ (equivalently same sizes). For every fixed $(i,j,z)$ there is a 1-1 matching between
\[\{(R_l,R_m,R_r): (R_l,R_m,R_r)\to(R'_l,R^*_m,R'_r),\ (e(R_l),e(R_r),z(R_l,R_m,R_r))=(i,j,z)\}\tag{5.35}\]
and the analogous set with $(R''_l,R^*_m,R''_r)$ (5.36). The matching fixes each middle $R_m$.
\begin{proof}Give a reversible map from (5.35) to (5.36). Take a triple in (5.35). By Remark5.22(2), the part of $R_l$ right of $S'_l$ is in $R^*_m$, hence disjoint from $R'_l$ and $R''_l$; symmetrically for the right side. Map
\[(R_l,R_m,R_r)\mapsto(\phi(R_l),R_m,\phi(R_r)),\]
where $\phi(R_l)$ replaces $R'_l$ by $R''_l$ in $R_l$, and symmetrically on the right. $R^*_m$, hence $R_m$, is unchanged. Since $R'_l,R''_l$ have the same size and the above parts are disjoint, this keeps $|V(R_l)|$. $R_l,\phi(R_l)$ have the same right distinguished set, the unique min-separator of both, so $e(R_l)=e(\phi(R_l))$; symmetrically on the right. Thus $(i,j)$ is unchanged. $z$ is unchanged since the changed part is disjoint from $Z(R_l,R_m,R_r)$. Reverse by switching the roles of the two ribbon pairs.\end{proof}
\paragraph{Lemma 5.25 (One round of factorization).} Let $L$ be (5.19) and $Q$ any $\binom{[n]}{\le d/2}\times\binom{[n]}{\le d/2}$ matrix with entries
\[Q(A,B)=\sum_{T_m:|V(T_m)\cup A\cup B|\le\tau}(\omega/n)^{|V(R_m)|}q(R_m)\chi_{T_m},\tag{5.37}\]
where $R_m=(A,B;T_m)$ and $q$ is symmetric w.r.t. shapes. Define $Q'$ and $E'_{\mathrm{negl}}$ so that
\[(LQL^\top)_{\mathrm{non\text{-}can}}=(LQ'L^\top)_{\mathrm{can}}+E'_{\mathrm{negl}}.\tag{5.38}\]
First set
\[Q'(A,B)=\sum_{T_m:|V(T_m)\cup A\cup B|\le\tau}(\omega/n)^{|V(R_m)|}q'(R_m)\chi_{T_m}.\tag{5.39}\]
Fix $R_m=(A,B;T_m)$ with $t=|V(R_m)|\le\tau$, $s=(|A|+|B|)/2$. For each generalized ribbon $R^*_m$ whose largest ribbon is $R_m$ and $|V(R^*_m)|\le\tau$, fix $(R'_l,R'_r)$ such that $(R'_l,R^*_m,R'_r)$ is the separating factorization of some triple with $|V(R'_l)|,|V(R'_r)|\le\tau$; exclude $R^*_m$ if none. Let
\begin{align*}q'(R_m)&=\sum_{\substack{R^*_m:\text{generalized on }(A,B)\\|V(R^*_m)|\le\tau,\ \text{largest ribbon }R_m}}(\omega/n)^{|I(R^*_m)|}q''(R^*_m),\\q''(R^*_m)&=\sum_{1\le z\le d/2}\sum_{\substack{\mathcal P=(R_l,R,R_r):\text{ side-inner-canonical}\\\mathcal P\to(R'_l,R^*_m,R'_r),\ z(\mathcal P)=z}}(\omega/n)^zq(R).\tag{5.40}\end{align*}
By Proposition5.24 this is independent of the fixed pair and remains symmetric w.r.t. shapes. Define $E'_{\mathrm{negl}}$ by (5.38). Conclusions: (1) w.p. $>1-n^{-9\log n}$,
\[\|E'_{\mathrm{negl}}\|\le\max\{q(A,B;T)\}\,n^{-\epsilon\tau}.\]
(2) If $C$ satisfies
\[\forall R_m\quad|q(R_m)|\le C(\omega/n^{1-\epsilon})^{s-p},\tag{5.41}\]
where $p$ is the maximum number of vertex-disjoint paths between $A,B$ in $R_m$ (equal to $s_{A,B}(T_m)$ by Menger), then
\[\forall R_m\quad|q'(R_m)|\le C(\omega/n^{1-\epsilon})^{s-p+1/3}.\]
\begin{proof}Compare the canonical and non-canonical parts as (0), then prove (1),(2).
\paragraph{(0).} For fixed $(I,J)$, $[LQL^\top]_{\mathrm{non\text{-}can}}(I,J)$ is
\[\sum_{\substack{(R_l,R_m,R_r):\text{side-inner-canonical, non-outer-canonical}\\\text{all three sizes}\le\tau}}(\omega/n)^{|V(R_l)|+|V(R_m)|+|V(R_r)|-|A|-|B|}q(R_m)\chi_{T_l\oplus T_m\oplus T_r},\tag{5.42}\]
where $(A,B)$ are the distinguished sets of $R_m$. Each triple has a unique separating factorization $(R'_l,R^*_m,R'_r)$.
\paragraph{First case: $|V(R^*_m)|\le\tau$.} The corresponding term in $(LQ'L^\top)_{\mathrm{can}}(I,J)$ is
\begin{multline*}(\omega/n)^{|V(R'_l)|+|V(R'_m)|+|V(R'_r)|-|S'_l|-|S'_r|}\\\cdot(\omega/n)^{z+|I(R^*_m)|}q(R'_m)\chi_{T'_l\oplus T^*_m\oplus T'_r},\tag{5.43}\end{multline*}
where $R'_m$ is the largest ribbon of $R^*_m$, $\chi_{T^*_m}$ means the character from $R'_m$, and $z\ge1$ is the original intersection size. Recall $T'_l\oplus T^*_m\oplus T'_r=T_l\oplus T_m\oplus T_r$ and
\begin{align*}|V(R_l)\cup V(R_m)\cup V(R_r)|&=|V(R'_l)|+|V(R^*_m)|+|V(R'_r)|-|S'_l|-|S'_r|\\&=|V(R_l)|+|V(R_m)|+|V(R_r)|-|A|-|B|-z.\end{align*}
Also $|V(R^*_m)|=|V(R'_m)|+|I(R^*_m)|$. Together the coefficient in (5.43) equals the one in (5.42) from $(R'_l,R^*_m,R'_r)$. Conversely, by $Q'$'s definition, (5.40) and Proposition5.24, each outer-canonical product corresponds uniquely to a side-inner-canonical triple of this case. Therefore the error collects precisely the next case.
\paragraph{Second case: $|V(R^*_m)|>\tau$.} We have
\begin{multline*}E'_{\mathrm{negl}}(I,J)=\sum_{\substack{(R_l,R_m,R_r):\text{side-inner-canonical, non-outer-canonical}\\\text{all three sizes}\le\tau,\ |V(R^*_m)|>\tau}}\\
(\omega/n)^{|V(R_l)|+|V(R_m)|+|V(R_r)|-|A|-|B|}\\
q(R_m)\chi_{T_l\oplus T_m\oplus T_r}.\tag{5.44}\end{multline*}
The obvious left/right distinguished-set condition $I,J$ is omitted.
\paragraph{(1).} For such a triple, putting $T=T_l\oplus T_m\oplus T_r$, the exponent is
\begin{align*}|V(R_l)|+|V(R_m)|+|V(R_r)|-|A|-|B|&=|V(R_l)\cup V(R_m)\cup V(R_r)|+z\\&=|V(T)\cup I\cup J|+|I(R^*_m)|+z.\end{align*}
Also $|I(R^*_m)|\le z+d/2$ from Lemma5.23.\footnote{In fact $|I(R^*_m)|\le z$ can be shown, but is not needed.} Fix $T$ and $a>\tau-|V(T)\cup I\cup J|$. To count triples contributing $(\omega/n)^{|V(T)\cup I\cup J|+a}\chi_T$ (temporarily ignore $q$), choose $I(R^*_m)$ of size $\le a/2+d/4$, since $a=|I(R^*_m)|+z\ge2|I(R^*_m)|-d/2$. Over fixed vertices the triple has $<3^{3\tau}2^{3\binom\tau2}$ choices. With $B_0=\max\{q(\cdot)\}$, its coefficient magnitude is less than
\begin{align*}&B_0(\omega/n)^{|V(T)\cup I\cup J|+a}n^{(a+d)/2}2^{2\tau^2}\\&=B_0(\omega/n^{1-2\epsilon})^{|V(T)\cup I\cup J|}n^{-2\epsilon|V(T)\cup I\cup J|}(\omega/\sqrt n)^an^{d/2}2^{2\tau^2}\\&\le B_0(n^{-1/2})^{|V(T)\cup I\cup J|}n^{-2\epsilon(|V(T)\cup I\cup J|+a)}n^{d/2}2^{2\tau^2}\\&\le B_0(n^{-1/2})^{|V(T)\cup I\cup J|}n^{-1.5\epsilon\tau}.\end{align*}
Use $\omega\le n^{1/2-4\epsilon}$, $|V(T)\cup I\cup J|+a>\tau$, $d<\epsilon\tau/10$, and $2^{2\tau}<n^{\epsilon/10}$. All appearing characters have $|V(T)|\le3\tau$. By Lemma4.2, for fixed $(I,J)$, w.p. $>1-n^{-10\log n}$,
\[|E'_{\mathrm{negl}}(I,J)|<\sum_{a=0}^{3\tau}B_0n^{-a/2}n^{-1.5\epsilon\tau}n^{a/2}n^{4\log\log n}2^{a^2}<n^{-1.4\epsilon\tau}.\]
Union-bound over $|\{(I,J)\}|<n^d$: w.p. $>1-n^{-9\log n}$, $\|E'_{\mathrm{negl}}\|<n^d n^{-1.4\epsilon\tau}<n^{-\epsilon\tau}$.
\paragraph{(2).} Fix $R_m$. By (5.40),
\begin{align*}q'(R_m)=\sum_{\substack{z,R^*_m:\text{largest ribbon }R_m}}(\omega/n)^{|I(R^*_m)|+z}\sum_{\substack{\mathcal P=(R_l,R,R_r):\text{side-inner-canonical}\\\mathcal P\to(R'_l,R^*_m,R'_r),\ z(\mathcal P)=z}}q(R).
\end{align*}
For fixed $R^*_m$ there are at most $8^{z\tau}<n^{\epsilon z}$ triples in the inner sum, since $(R'_l,R'_r)$ is fixed: after choosing each vertex's occurrence in the ribbons and $A,B\subseteq R^*_m$, only edges appearing more than once remain to assign. Each such edge has an endpoint in the fixed multiset $Z$ of size $\le z$. By Lemma5.23(2) and (5.41), the sum in absolute value is bounded by
\begin{align*}&n^{\epsilon z}(\omega/n)^{z+|I(R^*_m)|}|q(R)|\\&\le(\omega/n^{1-\epsilon})^{2(s'-s)+(p-p')+2|I(R^*_m)|}C(\omega/n^{1-\epsilon})^{s-p}\\&\le C(\omega/n)^{2|I(R^*_m)|}(\omega/n^{1-\epsilon})^{s'-p'+1/2}.\end{align*}
Here $(s,p)$ are the parameters of $R$ and $(s',p')$ those of $R_m$; the last step uses $s'-s\ge1/2$ from Lemma5.23(1). For fixed $i_0$ there are $<n^{i_0}$ choices of $R^*_m$ with $|I(R^*_m)|=i_0$, and $1\le z\le3\tau$. Hence
\begin{align*}|q'(R_m)|&\le3\tau\sum_{i_0=0}^{d/2}Cn^{i_0}(\omega/n)^{2i_0}(\omega/n^{1-\epsilon})^{s'-p'+1/2}\\&\le C(\omega/n^{1-\epsilon})^{s'-p'+1/3}.\end{align*}\end{proof}
Apply Lemma5.25 to the non-canonical term of (5.30), with $Q\leftarrow DQ_0D$: it becomes $[L(DQ_1D)L^\top]_{\mathrm{can}}+E'_{1;\mathrm{negl}}$. Repeat to obtain
\begin{align*}M'&=L[D(Q_0-Q_1+Q_2-\cdots\pm Q_d)D]L^\top\\&\quad-(E_{1;\mathrm{deg}}-\cdots\pm E_{1+d;\mathrm{deg}})+(E'_{1;\mathrm{negl}}+\cdots+E'_{d;\mathrm{negl}}).\tag{5.45}\end{align*}
\paragraph{Proposition 5.26 ([5],Claim6.15).} $Q_{d+1}=0$.
\begin{proof}Induct that every ribbon $R_m=(A,B;T_m)$ in $Q_k$ has $|A|+|B|\ge k$. The case $k=0$ is trivial. Every ribbon in $Q_{k+1}$ is the largest ribbon of a generalized middle in the separating factorization of a non-outer-canonical triple in $L(DQ_kD)L^\top$. The old middle has $|A|+|B|\ge k$, and Lemma5.23(1) gives $|A'|+|B'|\ge|A|+|B|+1\ge k+1$. For $k=d+1$ this is impossible with both distinguished sets in $\binom{[n]}{d/2}$.\end{proof}
\paragraph{Remark 5.27.} The final non-exact PSDness steps are omitted by the author, since similar arguments are given in the exact case, Section6. The selected target is the exact theorem; those complete Section6 steps follow below.
